\chapter{Configuración de entrenamiento de los modelos}
\label{anexo:configuracion-entrenamiento}

Este anexo reúne la configuración con la que se entrenaron los dos modelos de segmentación, con el detalle necesario para que un tercero pueda reproducir los resultados informados en el Capítulo 13. Se incluyen los parámetros de optimización, la estrategia de aumentación de datos y los controles de calidad aplicados al cierre del entrenamiento. % CROSS_REFERENCE_MAPPING_REQUIRED: Capítulo 13

% SOURCE_NUMBERING_ANOMALY_PRESERVED: F.1 literal in DOCX Heading2
\phantomsection
\section*{F.1 Parámetros de entrenamiento}
\addcontentsline{toc}{section}{F.1 Parámetros de entrenamiento}
\label{anexo:entrenamiento-f1-parametros}

\begin{table}[H]
\small
\centering
\caption{Configuración comparada del entrenamiento de ambos modelos.}
\label{tab:anexo-e-parametros-entrenamiento}
\begin{tabularx}{\textwidth}{|>{\raggedright\arraybackslash}p{0.25\textwidth}|>{\raggedright\arraybackslash}X|>{\raggedright\arraybackslash}X|}
\hline
\textbf{Parámetro} & \textbf{Modelo sagital} & \textbf{Modelo axial} \\
\hline
Conjunto de datos & SPIDER & Sudirman y Al-Kafri \\
\hline
Arquitectura & U-Net bidimensional, variante sagital & U-Net bidimensional, variante axial \\
\hline
Número de clases & 4 & 6 \\
\hline
Canales base & 16 & 16 \\
\hline
Dimensión de entrada & $256 \times 256$ & $256 \times 256$ \\
\hline
Semilla & 20260716 & 2026 \\
\hline
Tamaño de lote & 8 & 8 \\
\hline
Épocas máximas & 80 & 80 \\
\hline
Tasa de aprendizaje & $8 \times 10^{-4}$ & $8 \times 10^{-4}$ \\
\hline
Decaimiento de pesos & $1 \times 10^{-4}$ & $1 \times 10^{-4}$ \\
\hline
Paciencia de detención temprana & 12 & 12 \\
\hline
Paciencia del planificador & 4 & No aplicada \\
\hline
Recorte de gradiente & No aplicado & 1,0 \\
\hline
Precisión mixta & Sí & Sí \\
\hline
Función de pérdida & 0,5 $\times$ entropía cruzada ponderada + 0,5 $\times$ Dice sin fondo & Entropía cruzada ponderada + Dice suave sin fondo \\
\hline
Pesos de clase & Estimados del conjunto & Estimados del conjunto, con refuerzo sobre la clase de fondo de origen \\
\hline
Épocas completadas & 75, con detención temprana & Detención temprana \\
\hline
Época seleccionada & 63 & 17 \\
\hline
\end{tabularx}
\end{table}

La función de pérdida combina en ambos casos un término de entropía cruzada ponderada por clase con un término basado en el coeficiente Dice que excluye el fondo. La combinación responde a que la entropía cruzada por sí sola tiende a favorecer las clases mayoritarias, mientras que el término de Dice equilibra la contribución de las estructuras de menor tamaño.

% MIGRATION_PENDING: F.2 / Tabla E.2
% MIGRATION_PENDING: F.3 / Tabla E.3
% MIGRATION_PENDING: F.4
